\section{Introduction}
\label{sec:introduction}

Automated facial demographic attribute recognition plays a pivotal role across various intelligent application domains, including digital forensics systems, biometric access control, human-computer interaction, and personalized interactive services \cite{ref1}, \cite{ref2}. However, multiple studies have reported performance disparities across specific demographic subgroups, particularly females and individuals with darker skin tones \cite{ref3}, \cite{ref4}. Such disparities can be influenced by various factors, including data distribution imbalances and the limitations of feature representations in handling visual variations \cite{ref5}. The limitations of conventional systems in coping with such visual variations demonstrate that fair and accurate demographic recognition still faces tangible practical challenges, thereby necessitating novel representation approaches capable of modeling the interaction between racial and gender demographic attributes in an integrated and comprehensive manner.

While some facial attribute recognition approaches model attributes separately, multi-task learning enables multiple attributes to be learned simultaneously \cite{ref6}. Modeling attributes in isolation can restrict the analysis of groups formed by attribute combinations, such as race and gender \cite{ref7}, \cite{ref8}. Consequently, aggregate metrics potentially fail to capture performance disparities across intersectional demographic subgroups \cite{ref9}. Controlled evaluation across demographic subgroups with balanced representation can provide a more consistent and objective comparative foundation for measuring inter-group performance variations \cite{ref9}. Intersectional modeling also faces representational challenges because facial characteristics and expressions can vary across demographic groups, thereby necessitating visual representations capable of preserving discriminative information while remaining robust against such variations \cite{ref10}.

The methodological evolution of facial attribute recognition has shifted from handcrafted features toward deep learning paradigms, including CNN and ViT. In the era of Convolutional Neural Networks (CNN), various studies explored demographic feature extraction by leveraging specific facial regions such as the mid-face area, multi-scale multi-task frameworks for simultaneous age and gender estimation, and deep CNN models for ethnicity recognition \cite{ref11}, \cite{ref12}, \cite{ref13}. ViT is capable of modeling global inter-patch relationships via Multi-Head Self-Attention (MHSA), thereby potentially yielding more contextual visual representations, while the utilization of synthetic generative representations has also been developed to mitigate demographic disparities \cite{ref14}, \cite{ref15}. In addition, the integration of parallel attention mechanisms, the application of ViT for facial gender classification, and the exploration of multi-axis transformer architectures have further enriched the diversity of visual representation modeling \cite{ref16}, \cite{ref17}, \cite{ref18}. Furthermore, exploration into combining cross-domain ViT representations has emerged through the integration of dual facial and emotional features to enhance race and gender attribute discrimination \cite{ref19}, \cite{ref20}. Nevertheless, the fusion of biometric, expression, and aging representations into a unified representation remains relatively limited.

Despite the continuous advancements in ViT, several research gaps directly relevant to intersectional demographic classification remain. First, the majority of prior studies rely on single-domain facial biometric representations, while the simultaneous integration of affective information for intersectional demographic classification remains limited \cite{ref11}, \cite{ref17}, \cite{ref18}. Second, demographic recognition frequently exhibits high sensitivity to age-related facial morphological changes because the employed feature extraction does not incorporate biological aging representations \cite{ref12}, \cite{ref17}. Third, comparative studies regarding the influence of classifiers, feature scaling, and dimensionality reduction on transformer representations remain limited \cite{ref13}, \cite{ref19}, \cite{ref20}. To bridge the limitations of single-domain representations and evaluate these downstream classifier decision boundaries, this study proposes an integration of three ViT representation domains combined with structured classical classifier pipeline optimization.

To address these research gaps, this study proposes a Tri-Domain ViT feature fusion framework that combines three task-associated latent representations for intersectional demographic classification from facial images. The proposed framework integrates three complementary representation domains, including facial biometric representations, affective expressions, and aging characteristics, into a unified representation via feature fusion \cite{ref21}, \cite{ref22}. Latent feature extraction is conducted offline from frozen pre-trained ViT backbone models, allowing downstream classification to be performed without fine-tuning the backbone \cite{ref23}. The resulting feature vectors are subsequently evaluated through single-domain, dual-domain, and tri-domain ablation schemes using four classical machine learning classifiers, namely RF, GNB, LR, and SVM. All downstream classifiers are evaluated using Stratified Cross-Validation within a pipeline that integrates feature scaling and PCA dimensionality reduction, with each transformation learned strictly on the training data of each fold to prevent information leakage \cite{ref24}.

The main scientific contributions of this study are summarized into four primary points as follows:
\begin{enumerate}
\item A Tri-Domain ViT feature fusion framework integrating face-associated biometric representations, expression-related representations, and age-associated facial representations into a unified latent feature vector for intersectional demographic classification.
\item An empirical comparative benchmark across multiple feature ablation schemes and classical machine learning classifiers optimized via Stratified Cross-Validation hyperparameter tuning, examining decision boundary behavior across linear, probabilistic, ensemble, and kernel-based models.
\item Competitive classification performance on the DemogPairs dataset, with the proposed tri-domain representation outperforming the evaluated single-domain and dual-domain configurations for the majority of classifiers.
\item A granular intersectional subgroup performance analysis, evaluating subgroup-level classification performance across six demographic subgroups, providing an in-depth characterization of model behaviors across racial and gender intersections while systematically revealing the consistency of multi-domain feature representations on independent evaluation cohorts.
\end{enumerate}

The remainder of this article is organized as follows: Section~\ref{sec:related_works} discusses prior research on facial demographic recognition and transformer architectures. Section~\ref{sec:materials_and_methods} details the DemogPairs dataset specifications, multi-domain ViT representation extraction, and the optimized classification pipeline. Section~\ref{sec:results_and_discussion} presents the experimental evaluation results, feature ablation studies, and intersectional subgroup performance analysis. Finally, Section~\ref{sec:conclusion} summarizes the primary research conclusions, methodological limitations, and recommendations for future research directions.
